\section{Rotation of the fair-signed height at \texorpdfstring{$q=4$}{q=4}}
\label{sec:q4-rotation-repair}

The endpoint argument needs rotational invariance of height correlations,
with the loops counted with their multiplicities.  The loop-set distance
in the statement of the FK rotation theorem does not by itself retain
those multiplicities.  We obtain the required assertion directly from
the homotopy coupling used in its proof
\cite[Sections~3.5--3.6 and~4.2]{DCKKMO2026Rotation}.
Only qualitative bulk RSW estimates are used in this transfer.

For $k\ge2$, let
\[
 \mathcal D_k=
 \left\{\boldsymbol u=(u_1,u'_1,\ldots,u_k,u'_k):
 \{u_i,u'_i\}\cap\{u_j,u'_j\}=\varnothing\quad(i\ne j)\right\}.
\]
Coincidence within one pair is allowed.  Give the critical $q=4$ FK
loops independent fair signs and let $h_\delta$ be the resulting height,
with jumps of size one and an immaterial additive constant.  At a
continuum point use any fixed face representative within $O(\delta)$.
Put
\begin{equation}
 \Phi_{k,\delta}(\boldsymbol u)
 =\mathbb E\prod_{i=1}^k
       \bigl(h_\delta(u'_i)-h_\delta(u_i)\bigr).
 \label{eq:rr-Phi}
\end{equation}

\begin{theorem}[Isometry invariance of the endpoint correlations]
\label{thm:rr-rotation}
For every even $k\ge2$ and every compact $K\subset\mathcal D_k$,
\begin{equation}
 \lim_{\delta\downarrow0}
 \sup_{\boldsymbol u\in K}\sup_{I\in\operatorname{Isom}(\R^2)}
 \left|\Phi_{k,\delta}(\boldsymbol u)
             -\Phi_{k,\delta}(I\boldsymbol u)\right|=0.
 \label{eq:rr-target}
\end{equation}
\end{theorem}

The exact plane correspondence is recorded first, including the
passage from torus loops to the plane.  This avoids any need for an
unpublished general-parameter correspondence.

\begin{lemma}[The endpoint plane correspondence]
\label{lem:rr-bkw}
After the fixed similarity identifying the medial square grid with
$\Z^2$, the fair unit-jump height of critical plane $\FK_4$ has the
same gradient law as the zero-slope six-vertex height with
$\mathbf a=\mathbf b=1$ and $\mathbf c=2$.
\end{lemma}

\begin{proof}
On an even square torus, give equal weight to each pair consisting
of noncrossing local switch pairings and an orientation of each
resulting closed loop.  Erasing orientations gives weight $2^\ell$
to a configuration with $\ell$ loops.  Erasing pairings gives an
ice-rule arrow configuration: each of the four $a$- and $b$-type
vertices admits one compatible pairing, and each of the two
$c$-type vertices admits two.  Thus the arrow marginal has weights
$(1,1,2)$, and the loop orientations, conditionally on the
unoriented loops, are independent and fair.  This is the endpoint
specialization of the local calculation in
\cite[Section~2, Equations~(2.1)--(2.3)]{Lis2021}.

The torus weight $2^\ell$ need not equal the torus FK weight.
Precisely the needed plane passage is supplied by
\cite[Lemma~3.1]{Lis2021}: its unoriented loop law converges locally
to critical plane $\FK_4$.  Fix a finite set of observed medial
edges.  Bulk annular RSW implies that every plane FK loop meeting
this set is finite almost surely.  For any $\varepsilon>0$, choose
a box containing all these loops with probability at least
$1-\varepsilon$.  The event that the observed loops close inside
this box is determined by its finitely many pairings.  On that
event, local convergence transfers their full loop identities,
and hence their independent fair orientations.  Letting the box
grow proves convergence of the local arrow law to independently
oriented plane FK loops.

On the other hand, the unrestricted six-vertex torus laws have the
zero-slope plane limit used in
\cite[Theorem~2.2 and the paragraph following it]{SixVertexGFF2026}.
That paragraph identifies the unrestricted torus limit with the
slope-zero law defined through balanced cylinders.  This is a
lattice Gibbs-law fact preceding the GFF
argument.  The two descriptions of the arrow limit agree, and
arrows determine unit height differences.  Additive constants do
not affect the gradient law.  This proves the lemma.
\end{proof}

Consequently Theorem~\ref{thm:rr-rotation} supplies exactly the
$\mathbf c=2$ case of the symmetry input
\cite[Theorem~4.1]{SixVertexGFF2026}.  Its proof below does not use
any GFF limit or any sharp four-arm estimate.

\subsection{Even moments retain only macroscopic loops}

Smooth the discrete contours locally into disjoint Jordan loops.  This
changes their location by $O(\delta)$.  For such a loop $L$ set
\[
 D_i(L)=\mathbf1_{\operatorname{int}L}(u'_i)
                   -\mathbf1_{\operatorname{int}L}(u_i).
\]
A reference orientation may multiply all the $D_i(L)$ by one sign
depending on $L$; the fair sign absorbs that choice.  Let
$\Pi_{\mathrm{even}}(k)$ denote the partitions of $\{1,\ldots,k\}$
whose blocks all have even size.  Conditional expectation over the
signs gives the exact identity
\begin{equation}
 P_k(\mathcal L;\boldsymbol u)
 :=\mathbb E\left[\prod_i
       \bigl(h_\delta(u'_i)-h_\delta(u_i)\bigr)
                         \,\middle|\,\mathcal L\right]
 =\sum_{\pi\in\Pi_{\mathrm{even}}(k)}
   \ \sum_{\substack{f:\pi\to\mathcal L\\f\ \mathrm{injective}}}
          \prod_{B\in\pi}\prod_{i\in B}D_i(f(B)).
 \label{eq:rr-even-polynomial}
\end{equation}
At positive mesh only finitely many loops separate any prescribed
pair: each such loop crosses a fixed finite lattice path between its
endpoints.  Hence the expansion is a finite identity.  In particular,
there is no limiting interchange in its derivation.

Call a loop relevant if at least two coordinates of
$(D_1(L),\ldots,D_k(L))$ are nonzero.  Every loop in a nonzero summand
of \eqref{eq:rr-even-polynomial} is relevant.  If endpoints belonging
to distinct slots are separated by $d>0$, a relevant loop surrounds
one endpoint from each of two distinct slots.  Therefore
\begin{equation}
 \diam L\ge d.
 \label{eq:rr-macro}
\end{equation}
Here we used that the closure of a Jordan domain is contained in the
convex hull of its boundary.  Also, a separating loop intersects the
segment joining the separated endpoints.  Thus, if the endpoints lie
in $B(0,M)$, every relevant loop meets that disk.  For rounded
continuum endpoints, replace $d$ by $d-O(\delta)$.

We record the quantitative estimates needed below.  They hold also
for the two-angle mixed isoradial lattices in the rotation
construction.  For any fixed choice of the two angles, their mesh
geometry has bounded distortion and their bulk conditional RSW
bounds are available from
\cite[Theorem~2.2, Remark~2.3 and Proposition~2.12]{DCKKMO2026Rotation}.

\begin{lemma}[Counts, escape, and deterministic-point avoidance]
\label{lem:rr-estimates}
Let $N(a,M)$ be the number of loops of diameter at least $a$ meeting
$B(0,M)$.  For $\delta$ sufficiently smaller than $a$, and every
fixed $p<\infty$,
\begin{equation}
 \|N(a,M)\|_{L^p}\le C_p(1+M/a)^2.
 \label{eq:rr-count}
\end{equation}
There are $C,\gamma>0$ such that, with harmless fixed changes of the
radii,
\begin{align}
 \mathbb P\{\exists L:L\cap B(0,M)\ne\varnothing,
                            L\not\subset B(0,R)\}
      &\le C(M/R)^\gamma, && R\ge10M,
       \label{eq:rr-escape}\\
 \mathbb P\{\exists L:\diam L\ge a,
                       \dist(z,L)\le r\}
      &\le C\bigl((r+C\delta)/a\bigr)^\gamma,
          &&0<r<a/100.
       \label{eq:rr-point}
\end{align}
The constants are independent of the deterministic center $z$.
\end{lemma}

\begin{proof}
Bulk RSW in a sequence of buffered annuli gives a polynomial bound
for an actual primal or dual arm between their extreme rims: each
annulus has a uniformly positive conditional probability of a
blocking circuit of the opposite color.

We also need the standard elementary count consequence of this
argument.  The number of distinct actual clusters of one color
crossing a fixed-aspect annulus has a geometric tail.  Explore
complete actual components from the inner rim, revealing their
closed exiting edges.  Remove every explored component, including
noncrossing ones, until the next crossing component is found.
At this stopping time, extend the surviving configuration by closed
edges on the deleted graph.  Its FK law is dominated by the law in
the original annulus with all vertices on both original rims
wired together.  Removed bulk
components do not wire surviving bulk vertices.  A dual circuit in
a buffered middle annulus of this auxiliary residual configuration
blocks another actual primal crossing, and has probability bounded
below even under the maximal primal boundary condition.  Repeating the exploration bounds the chance
of each additional crossing by a fixed number below one.  Apply
the same argument with the colors exchanged.

Each loop is the unique exterior contour of its interior-side
actual cluster, primal or dual.  Thus distinct loops give distinct
colored clusters.  Cover $B(0,M)$ by $O((1+M/a)^2)$ balls of radius
$a/100$.  A loop of diameter at least $a$ meeting one of these balls
has an interior-side bank crossing an annulus about its center
whose two radii are fixed multiples of $a$.  Different global
clusters give different actual components in that annulus.
The preceding geometric tail and Minkowski's inequality prove
\eqref{eq:rr-count}.

For \eqref{eq:rr-escape}, a bank of an escaping loop gives an
actual arm from scale $M$ to scale $R$.  For \eqref{eq:rr-point}, a
bank of the indicated loop gives an arm from scale $r+C\delta$
about $z$ to a fixed fraction of $a$.  The polynomial arm bound
proves both assertions.
\end{proof}

Fix a compact $K\subset\mathcal D_k$, write $d_K>0$ for its minimum
cross-slot endpoint distance, and choose $M$ containing every
endpoint in $K$.  If $k=2m$, an even partition has at most $m$
blocks, so \eqref{eq:rr-even-polynomial} implies
\begin{equation}
 |P_{2m}(\mathcal L;\boldsymbol u)|
       \le C_m N(d_K/2,M+1)^m,
 \qquad \boldsymbol u\in K,
 \label{eq:rr-UI}
\end{equation}
for sufficiently small mesh.  All fixed moments of the right-hand
side are uniformly bounded by Lemma~\ref{lem:rr-estimates}.  These
are bounds for the \emph{sign-averaged polynomial}; no uniform
absolute-moment assertion about a product of pointwise height
differences is being made.

Two consequences will be useful.  First, removing loops not
contained in $B(0,R)$ changes the expectation of the polynomial by
at most $C_K(M/R)^{\gamma/2}$, by
\eqref{eq:rr-escape}, \eqref{eq:rr-UI}, and Cauchy--Schwarz.
Second, if corresponding endpoints of two tuples in a compact
subset of $\mathcal D_k$ move by at most $r$, evaluate both
polynomials in the same loop sample.  Off the events in
\eqref{eq:rr-point} at those endpoints, every relevant loop has
unchanged endpoint indicators.  Hence
\begin{equation}
 |\Phi_{k,\delta}(\boldsymbol u)
                   -\Phi_{k,\delta}(\boldsymbol v)|
       \le C_K\bigl((r+C\delta)/d_K\bigr)^{\gamma/2}.
 \label{eq:rr-continuity}
\end{equation}
This estimate is uniform in the absolute location of the tuple.
It covers within-slot coincidences and changes of the rounding
convention.  The supremum over $K$ is a supremum of deterministic
tuple estimates; we do not require one avoidance event to hold
simultaneously at all points occurring in $K$.

\subsection{What the published coupling preserves}

Let $\mathbb L(\alpha)$ be the rectangular isoradial lattice in
\cite[Section~1.4]{DCKKMO2026Rotation};
$\mathbb L(\pi/2)$ is a fixed rotation and dilation of the square
lattice.  The next lemma extracts the precise part of its coupling
that we use.  Its assertion is about loop multiplicities, not merely
about proximity of loop sets.

\begin{lemma}[Homotopy coupling with dense marked clusters]
\label{lem:rr-nails}
Fix $\alpha\in(0,\pi)$.  As the mesh tends to zero, the construction
of \cite[Section~3.6]{DCKKMO2026Rotation} gives a coupling of two
auxiliary loop samples.  Their restrictions to a deterministic
window exhausting the plane have, unconditionally, the respective
homogeneous laws on $\delta\mathbb L(\pi/2)$ and
$\delta\mathbb L(\alpha)$.

Outside an event of probability tending to zero, there are
correspondingly indexed, pairwise disjoint marked
clusters, called nails, whose mesh tends to zero and whose window
exhausts the plane.  Corresponding nails have diameter tending to
zero and lie at distance tending to zero from the same grid point.
There is a bijection $\psi$ between all loops surrounding at least
two but not all nails in the two samples.  It preserves their
oriented reduced crossing words relative to the nail grid.
\end{lemma}

\begin{proof}
This is the construction preceding the metric conclusion, with its
parameters retained.  In the notation of
\cite[Definition~3.8 and Section~3.6, pp.~49--50]{DCKKMO2026Rotation},
choose $\eta=N^{-b}$ and $t=N^{-b^2}$, for $b>0$ sufficiently small.
On an event of probability $1-O(N^{-c})$, the selected nail
clusters survive the whole extrema-resampling and track-exchange
chain.  The same construction preserves the bijection of
non-trivial loops and their reduced words at each step.
This covers both colors of contour.

Initially, the nail grid spacing is $t^{1/2}N$, nail diameters
are $O(tN)$, and the nail window has size comparable to
$t^{1/4}N$.  The terminal nail positions differ by these errors
from their images under $M_{\pi/2,\alpha}$.  The independent FK
argument in \cite[Section~4.2, proof of Theorem~1.7]{DCKKMO2026Rotation}
proves $M_{\pi/2,\alpha}=\mathrm{id}$.  We may therefore use the
already constructed coupling with corresponding nails at the same
limiting positions.

For clarity, these scales can be arranged for every small physical
mesh $\delta$.  Set $a=b^2$ and choose
$N=2\lfloor\tfrac12\delta^{-1/(1-3a/8)}\rfloor$, so $N$
is even as required by the track construction.  Then the physical grid
spacing, nail diameter, and nail-window size are respectively of
orders
\[
 N^{-a/8},\qquad N^{-5a/8},\qquad N^{a/8}.
\]
The homogeneous observation window has scale $\delta N$, which
also diverges.  Center the nail window by translating each
endpoint sample by a period vector of its own homogeneous
lattice nearest the same centering vector.  The two chosen
translations differ by $O(\delta)$; this only adds a vanishing
error to corresponding nail positions.  The endpoint restrictions
have exactly the homogeneous laws, by
\cite[Proposition~2.11]{DCKKMO2026Rotation}.  After centering, both
the nail window and the homogeneous observation windows exhaust
the plane.  The probability $O(N^{-c})$ tends to zero.  This proves
the stated extraction; no conclusion about curve-set convergence
is used to recover the bijection.
\end{proof}

Here is the topological information in a crossing word that will
matter.  Orient each loop by the reference convention and let
$s_L\in\{-1,1\}$ be its winding number at an interior point.
For any two nails, join chosen points in them by a path of nail-grid
edges, with transverse perturbations if necessary.  The algebraic
intersection number of $L$ with this path equals the difference of
its winding numbers at the two nail points.  It is read from the
oriented crossing word.  Cyclic permutation and deletion of an
adjacent inverse pair do not change it.  Consequently, if $n,n'$
and $\widetilde n,\widetilde n'$ are corresponding nail points,
\begin{equation}
 s_L\bigl(\mathbf1_{\operatorname{int}L}(n')
                  -\mathbf1_{\operatorname{int}L}(n)\bigr)
 =s_{\psi(L)}
       \bigl(\mathbf1_{\operatorname{int}\psi(L)}(\widetilde n')
                  -\mathbf1_{\operatorname{int}\psi(L)}(\widetilde n)\bigr).
 \label{eq:rr-winding-difference}
\end{equation}
We need only these differences.  In particular, no inference of an
absolute inside set from a finite grid word is required.

\begin{figure}[t]
\centering
\begin{tikzpicture}[scale=.83,>=Stealth]
 \begin{scope}
  \draw[thick,blue!55!black] (0.8,0) ellipse (1.22 and 1.48);
  \node[blue!55!black] at (1.8,1.5) {$L$};
  \foreach \y/\j in {.7/1,-.7/2}{
   \fill (-1.25,\y) circle (1.6pt);
   \fill (.8,\y) circle (1.6pt);
   \node[anchor=east] at (-1.34,\y) {$u_{\j}$};
   \node[anchor=west] at (.88,\y) {$u'_{\j}$};
   \draw[gray,densely dashed] (-1.25,\y)--(-1.1,\y+.15);
   \draw[gray,densely dashed] (.8,\y)--(.65,\y+.15);
   \draw[red!70!black,fill=white,semithick] (-1.1,\y+.15) rectangle ++(.12,.12);
   \draw[red!70!black,fill=white,semithick] (.65,\y+.15) rectangle ++(.12,.12);
  }
  \node at (.35,-1.95) {two separated slots};
 \end{scope}
 \draw[->,thick] (2.4,0)--(4.1,0)
        node[midway,above] {$\psi$};
 \begin{scope}[xshift=6cm]
  \draw[thick,blue!55!black]
      (-.35,-1.1) .. controls (-.9,-.1) and (-.25,1.5) .. (.9,1.4)
      .. controls (2,1.2) and (2,-1.4) .. (.75,-1.5)
      .. controls (.1,-1.7) and (-.15,-1.4) .. (-.35,-1.1);
  \node[blue!55!black] at (1.8,1.65) {$\psi(L)$};
  \foreach \y/\j in {.7/1,-.7/2}{
   \fill (-1.25,\y) circle (1.6pt);
   \fill (.8,\y) circle (1.6pt);
   \draw[gray,densely dashed] (-1.25,\y)--(-1.1,\y+.15);
   \draw[gray,densely dashed] (.8,\y)--(.65,\y+.15);
   \draw[red!70!black,fill=white,semithick] (-1.1,\y+.15) rectangle ++(.12,.12);
   \draw[red!70!black,fill=white,semithick] (.65,\y+.15) rectangle ++(.12,.12);
  }
  \node at (.35,-1.95) {$D(\psi(L))=\pm D(L)$};
 \end{scope}
\end{tikzpicture}
\caption{Nearby nails (squares) detect the separations of endpoint
pairs (dots).  The homotopy bijection preserves oriented winding
differences.  A single sign change of a loop's whole separation
vector cancels when that loop occurs an even number of times.}
\label{fig:rr-nails}
\end{figure}

\subsection{Transfer of the correlation polynomial}

Fix $\alpha\in(0,\pi)$ and compact $K\subset\mathcal D_k$.
Choose $0<r<d_K/1000$.  Apply Lemma~\ref{lem:rr-nails} and, for
each endpoint, select an indexed nail close to it.  For all small
enough meshes both corresponding nails lie within $r/2$ of that
endpoint.  Nails selected for endpoints in distinct slots are
distinct.  In both samples discard the event that a loop of
diameter at least $d_K/4$ comes within $r$ of an endpoint.
Lemma~\ref{lem:rr-estimates} bounds the probability of these
exceptions by $C_K(r/d_K)^\gamma+o(1)$, uniformly for the
individual tuples in $K$.

On the remaining event, a relevant loop $L$ surrounds the nails
near one endpoint of each of two distinct slots, and excludes a
nail near the other endpoint of a separated slot.  It is therefore
non-trivial in the exact sense of Lemma~\ref{lem:rr-nails}, so
$\psi(L)$ is defined.  For those two slots the winding differences
at the corresponding nails are nonzero by
\eqref{eq:rr-winding-difference}.  Thus $\psi(L)$ separates each
of the two nail pairs.  In particular, it encloses one nail from
each of two distinct slots, whose distance is at least
$d_K-O(r+\delta)$.  This proves that $\psi(L)$ itself has diameter
at least $d_K/4$.  Its endpoint-avoidance event is consequently
applicable; smallness of the partner has not been assumed away.

Endpoint indicators and nearby-nail indicators agree for both
loops.  Equation~\eqref{eq:rr-winding-difference} now gives,
for \emph{every} slot,
\begin{equation}
 D_i(\psi(L))=\frac{s_L}{s_{\psi(L)}}D_i(L).
 \label{eq:rr-signature}
\end{equation}
The inverse bijection gives the same statement starting with a
relevant loop in the second sample.  Hence it is a bijection
between all relevant loops, preserving their entire separation
vectors up to one sign per loop.  For a block $B$ in an even
partition,
\[
 \prod_{i\in B}D_i(\psi(L))
    =\left(\frac{s_L}{s_{\psi(L)}}\right)^{|B|}
                       \prod_{i\in B}D_i(L)
    =\prod_{i\in B}D_i(L).
\]
Transporting the injective labels in
\eqref{eq:rr-even-polynomial} therefore makes the two polynomials
\emph{exactly equal} on the good event.  This uses the actual
bijection, so distinct loops are never identified with one another.
Cauchy--Schwarz and \eqref{eq:rr-UI} bound the expectation error by
$C_K(r/d_K)^{\gamma/2}+o(1)$.

It remains to pass from the auxiliary mixed samples to the intended
homogeneous ones.  Fix $R>10M$ and couple their restrictions
identically in $B(0,R)$ using the local-law assertion of
Lemma~\ref{lem:rr-nails}; for small enough mesh this disk lies in
both homogeneous windows.  Off the escaping-loop event in either
sample, all relevant complete loops are identical.  Equations
\eqref{eq:rr-escape} and \eqref{eq:rr-UI} bound the expectation
error by $C_K(M/R)^{\gamma/2}$.  Let the mesh tend to zero,
then $r\downarrow0$ and $R\uparrow\infty$.  Writing
$\Phi^{\alpha}_{k,\delta}$ for the fair-signed correlation on
$\delta\mathbb L(\alpha)$, we obtain
\begin{equation}
 \lim_{\delta\downarrow0}\sup_{\boldsymbol u\in K}
 \left|\Phi^{\pi/2}_{k,\delta}(\boldsymbol u)
                   -\Phi^{\alpha}_{k,\delta}(\boldsymbol u)\right|=0.
 \label{eq:rr-isoradial-transfer}
\end{equation}
No diameter threshold needs to be a continuity point: all small
loops were removed by the even-moment identity itself.

\begin{proof}[Proof of Theorem~\ref{thm:rr-rotation}]
The rectangular model on $\mathbb L(\alpha)$ is invariant under
reflection $S_{\alpha/2}$ about $e^{i\alpha/2}\R$
\cite[Remark~1.8]{DCKKMO2026Rotation}.  Its fair-sign height
correlations inherit that symmetry.  If the chosen endpoint
rounding is not reflection equivariant, the reflected
representatives differ by $O(\delta)$, which is harmless by
\eqref{eq:rr-continuity}.  Apply
\eqref{eq:rr-isoradial-transfer} at the original and reflected
tuples.  The square model is consequently asymptotically invariant
under $S_{\alpha/2}$, uniformly on every compact subset of
$\mathcal D_k$.  Composing with its exact horizontal-reflection
symmetry yields invariance under rotation through $\alpha$.
Square symmetries cover the endpoints and the other half of the
circle.  This proves the assertion for each fixed rotation.

The estimate \eqref{eq:rr-continuity} upgrades this to uniformity
in the rotation angle.  Indeed, the orthogonal images of $K$ form
a compact subset of $\mathcal D_k$, and a finite angular net
approximates every rotation with endpoint displacement at most
$C_K$ times its angular mesh.  First use convergence at the
finitely many net angles, and then let the angular mesh decrease.
Reflections follow by composition with an exact square symmetry.
This argument requires no estimate uniform in a degenerating
isoradial angle.

Finally, any translation differs from a period translation of the
mesh-$\delta$ square lattice by a vector of length $O(\delta)$.
Exact period invariance and the location-uniform estimate
\eqref{eq:rr-continuity} handle this residual displacement,
uniformly over \emph{all} translation vectors.  Every plane
isometry is a translation followed by an orthogonal map.
The fixed rotation and dilation identifying
$\mathbb L(\pi/2)$ with $\Z^2$ merely change the mesh convention
and compact set.  These observations prove
\eqref{eq:rr-target} for the original square lattice.
\end{proof}


\subsection{The whole-plane field input}

\begin{corollary}[The endpoint whole-plane field limit]
\label{cor:rr-gff-input}
The fair unit-jump plane height of critical square-lattice $\FK_4$
converges on neutral smooth compactly supported tests to
$\sqrt{2/\pi}$ times the whole-plane GFF with Green kernel
$G(x,y)=-(2\pi)^{-1}\log|x-y|$.  In particular,
Proposition~\ref{q4field:wholeplane-input} holds.
\end{corollary}

\begin{proof}
Lemma~\ref{lem:rr-bkw} identifies the discrete gradient law with
the isotropic six-vertex law at $(\mathbf a,\mathbf b,\mathbf c)
=(1,1,2)$.  Theorem~\ref{thm:rr-rotation} gives precisely the
compact-$\mathcal D_k$, all-isometries symmetry statement used as
Theorem~4.1 in \cite{SixVertexGFF2026}, at this parameter.
We may therefore apply its isotropic Theorem~2.8, using the
remaining spectral, regularity, and free-energy results cited and
proved there.  The passage from correlations to smooth field
tests is part of its Section~8.  No anisotropic extension or
general-parameter loop-functional assertion is needed here.

At $\mathbf c=2$ one has $\Delta=-1$, so the variance multiplier
in that theorem is $2/\arccos(-1)=2/\pi$.  A fixed rotation and
dilation of the lattice does not change the neutral whole-plane
Green quadratic form: the dilation adds a constant to the
logarithmic kernel and each test has integral zero.  For every
real $\rho\in C_c^\infty(\R^2)$ with $\int\rho=0$, the convergence
therefore gives
\[
 \mathbb E e^{i\langle h_\delta,\rho\rangle}
 \longrightarrow
 \exp\left\{-\frac1\pi
       \iint \rho(x)G(x,y)\rho(y)\,dx\,dy\right\}.
\]
Apply the same assertion to every real linear combination of a
finite family of neutral tests to obtain joint convergence.
This is exactly Proposition~\ref{q4field:wholeplane-input}.
\end{proof}
